% ==============================================================================
% CONCLUSIONES INDIVIDUALES
% ==============================================================================
% Instrucciones de la plantilla:
% Formato: #CUENTA: CONCLUSIÓN
% Cada integrante del equipo debe redactar de forma personal, concreta y 
% coherente sus conclusiones individuales sobre la práctica y los conceptos 
% gráficos aprendidos.
% ==============================================================================

\section{Conclusiones Individuales}

\subsection*{\#321137135: Basilio Illescas Andrés}
% Redactar la conclusión individual del alumno aquí:
A través de las actividades implementadas, se logró afianzar el conocimiento sobre...

\subsection*{\#424065465: Cruz Macedo Samuel Santiago}
% Redactar la conclusión individual del alumno aquí:
El análisis de los experimentos realizados permitió identificar...

\subsection*{\#321258173: Medel Piña Ángel de Jesús}
% Redactar la conclusión individual del alumno aquí:
En esta práctica entendí que, en el modelo de Phong, la luz de varias fuentes simplemente se suma. En el laboratorio lo vimos con cuatro fuentes de color, y en la actividad lo llevé a diez, que es el máximo que acepta el \textit{shader}. Con diez fuentes apareció un problema que con una sola no había notado: la componente ambiental se suma una vez por cada luz, así que con la potencia de 60 la escena se saturaba a blanco y tuve que bajarla a 30. No esperaba que, aunque los colores de las diez fuentes suman casi blanco, la casa siguiera viéndose de colores distintos según la zona; la razón es que la componente difusa y la especular dependen de dónde está cada fuente. Al dejar encendida solo la luz roja vi que el terreno detrás de la casa seguía iluminado, lo que me dejó claro que este modelo no calcula sombras. Con la tabla de materiales comprobé que cambiar unos cuantos coeficientes basta para pasar de una esfera gris a una esmeralda o un rubí, y con el efecto Fresnel, que para materiales como el vidrio o el diamante hace falta tomar en cuenta la refracción, algo que Phong no modela.

\subsection*{\#321150503: Pérez Olvera Alexis Abraham}
% Borrador para revisión y personalización del alumno.
Al trabajar con el toroide comprendí que importar un modelo no es suficiente para mostrarlo correctamente: también es necesario conservar sus normales, sus coordenadas UV y la relación con el archivo de textura. El tablero de ajedrez me permitió identificar cómo se distribuye la imagen sobre la superficie y comprobar que aumentar su repetición no modifica la geometría. Al comparar la vista sin iluminación con Phong, pude distinguir el color propio de la textura de los cambios producidos por la luz y el brillo especular. La diferencia de escala al importar el FBX también me hizo ver la importancia de revisar las unidades antes de ajustar la cámara. Considero que el resultado constituye una base funcional para incorporar materiales distintos, aunque esa ampliación todavía requiere configurar y evaluar las seis apariencias solicitadas en el manual.

\subsection*{\#320224159: Segura Cedeño Luisa María}
% Redactar la conclusión individual del alumno aquí:
Como resultado del desarrollo de los algoritmos y pruebas gráficas...
